% GENERATED FILE. Edit canonical lessons, then run npm run generate:books.
% canonical-source-hash: fnv1a64:8d4bc2eca4dfc3e8
% canonical-lessons: RU-C523-syuzhet, RU-C523-personazh, RU-C523-geroinya, RU-C523-razvyazka, RU-C523-zachityvatsya

\chapter{Plot and Characters}
\label{ch:ru-plot-and-characters}
\hlchaptermodality{523}

\begin{chapteropening}
\textbf{By the end of this chapter:} \emph{I can talk about the story of a book or a film.}

Complete the last lesson of chapter 523: I can talk about the story of a book or a film.
\end{chapteropening}

\section[\texorpdfstring{\ru{сюжет} (syuzhét)}{сюжет (syuzhét)}]{\ru{сюжет} (syuzhét) — a plot}

\label{lesson:RU-C523-syuzhet}



\begin{quote}
Before the new one: say the Russian for a statue, then the Russian for to inspire.
\end{quote}

\subsection*{You'll want to know: \ru{сюжет}}
\textbf{\ru{сюжет}} (\emph{syuzhét}) — \textquotedblleft{}a plot\textquotedblright{}.

\begin{quote}
\ru{У} \ru{этого} \ru{романа} \ru{интересный} \ru{сюжет}.
\end{quote}

\begin{quote}
\emph{U étogo romána interésnyy syuzhét.}
\end{quote}

\begin{quote}
This novel has an interesting plot.
\end{quote}

Masculine.

\begin{cousinweb}
\emph{\ru{сюжет}} was borrowed from French \emph{sujet} 'subject', which comes from Latin \emph{subiectum} (EtymDB), from \emph{subicere} 'to place under'.
\end{cousinweb}

\begin{grammarlens}[title={What you've built}]
One more word for telling, explaining and describing at B1.
\end{grammarlens}

\subsection*{Your turn}
\begin{itemize}
\raggedright
  \item \emph{Say it:} \emph{syuzhét}
  \item \emph{Say it:} \emph{syuzhét}, once more
  \item \emph{Say it:} a sentence of your own with \emph{syuzhét}
  \item \emph{Recall:} say the Russian for a statue, then the Russian for to inspire, then say \emph{syuzhét} again
\end{itemize}

\begin{tcolorbox}[breakable,colback=teal!4,colframe=teal!35!black,title={Before you move on}]
What does \ru{сюжет} mean? (\textquotedblleft{}A plot\textquotedblright{}.) Say it once more.
\end{tcolorbox}
\section[\texorpdfstring{\ru{персонаж} (personázh)}{персонаж (personázh)}]{\ru{персонаж} (personázh) — a character (in a story)}

\label{lesson:RU-C523-personazh}



\begin{quote}
Before the new one: say the Russian for to inspire, then the Russian for a plot.
\end{quote}

\subsection*{You'll want to know: \ru{персонаж}}
\textbf{\ru{персонаж}} (\emph{personázh}) — \textquotedblleft{}a character (in a story)\textquotedblright{}.

\begin{quote}
\ru{Мой} \ru{любимый} \ru{персонаж} — \ru{старый} \ru{капитан}.
\end{quote}

\begin{quote}
\emph{Moy lyubímyy personázh — stáryy kapitán.}
\end{quote}

\begin{quote}
My favourite character is the old captain.
\end{quote}

Masculine.

\begin{cousinweb}
\emph{\ru{персонаж}} is usually traced to French \emph{personnage}, built on \emph{personne}, from Latin \emph{persōna} 'a mask, a part in a play'; EtymDB has no entry.
\end{cousinweb}

\begin{grammarlens}[title={What you've built}]
One more word for telling, explaining and describing at B1.
\end{grammarlens}

\subsection*{Your turn}
\begin{itemize}
\raggedright
  \item \emph{Say it:} \emph{personázh}
  \item \emph{Say it:} \emph{personázh}, once more
  \item \emph{Say it:} a sentence of your own with \emph{personázh}
  \item \emph{Recall:} say the Russian for to inspire, then the Russian for a plot, then say \emph{personázh} again
\end{itemize}

\begin{tcolorbox}[breakable,colback=teal!4,colframe=teal!35!black,title={Before you move on}]
What does \ru{персонаж} mean? (\textquotedblleft{}A character (in a story)\textquotedblright{}.) Say it once more.
\end{tcolorbox}
\section[\texorpdfstring{\ru{героиня} (geroínya)}{героиня (geroínya)}]{\ru{героиня} (geroínya) — a heroine}

\label{lesson:RU-C523-geroinya}



\begin{quote}
Before the new one: say the Russian for a plot, then the Russian for a character.
\end{quote}

\subsection*{You'll want to know: \ru{героиня}}
\textbf{\ru{героиня}} (\emph{geroínya}) — \textquotedblleft{}a heroine\textquotedblright{}.

\begin{quote}
\ru{Героиня} \ru{романа} \ru{уезжает} \ru{в} \ru{другой} \ru{город}.
\end{quote}

\begin{quote}
\emph{Geroínya romána uyezzháyet v drugóy górod.}
\end{quote}

\begin{quote}
The heroine of the novel moves to another town.
\end{quote}

Feminine.

\begin{cousinweb}
\emph{\ru{героиня}} was borrowed from Greek \emph{hērōínē} 'a heroine' (EtymDB), the feminine of \emph{h\'{\={e}}rōs} 'hero'.
\end{cousinweb}

\begin{grammarlens}[title={What you've built}]
One more word for telling, explaining and describing at B1.
\end{grammarlens}

\subsection*{Your turn}
\begin{itemize}
\raggedright
  \item \emph{Say it:} \emph{geroínya}
  \item \emph{Say it:} \emph{geroínya}, once more
  \item \emph{Say it:} a sentence of your own with \emph{geroínya}
  \item \emph{Recall:} say the Russian for a plot, then the Russian for a character, then say \emph{geroínya} again
\end{itemize}

\begin{tcolorbox}[breakable,colback=teal!4,colframe=teal!35!black,title={Before you move on}]
What does \ru{героиня} mean? (\textquotedblleft{}A heroine\textquotedblright{}.) Say it once more.
\end{tcolorbox}
\section[\texorpdfstring{\ru{развязка} (razvyázka)}{развязка (razvyázka)}]{\ru{развязка} (razvyázka) — an ending, a denouement}

\label{lesson:RU-C523-razvyazka}



\begin{quote}
Before the new one: say the Russian for a character, then the Russian for a heroine.
\end{quote}

\subsection*{You'll want to know: \ru{развязка}}
\textbf{\ru{развязка}} (\emph{razvyázka}) — \textquotedblleft{}an ending, a denouement\textquotedblright{}.

\begin{quote}
\ru{Развязка} \ru{фильма} \ru{меня} \ru{удивила}.
\end{quote}

\begin{quote}
\emph{Razvyázka fílma menyá udivíla.}
\end{quote}

\begin{quote}
The ending of the film surprised me.
\end{quote}

Feminine.

\begin{cousinweb}
\emph{\ru{развязка}} is built on \emph{\ru{развязать}} 'to untie', from \emph{\ru{вязать}} 'to tie', which is inherited from Proto-Slavic \emph{vęzati} (EtymDB): the ending unties the knot of the story.
\end{cousinweb}

\begin{grammarlens}[title={What you've built}]
One more word for telling, explaining and describing at B1.
\end{grammarlens}

\subsection*{Your turn}
\begin{itemize}
\raggedright
  \item \emph{Say it:} \emph{razvyázka}
  \item \emph{Say it:} \emph{razvyázka}, once more
  \item \emph{Say it:} a sentence of your own with \emph{razvyázka}
  \item \emph{Recall:} say the Russian for a character, then the Russian for a heroine, then say \emph{razvyázka} again
\end{itemize}

\begin{tcolorbox}[breakable,colback=teal!4,colframe=teal!35!black,title={Before you move on}]
What does \ru{развязка} mean? (\textquotedblleft{}An ending, a denouement\textquotedblright{}.) Say it once more.
\end{tcolorbox}
\section[\texorpdfstring{\ru{зачитываться} (zachítyvatsya)}{зачитываться (zachítyvatsya)}]{\ru{зачитываться} (zachítyvatsya) — to get engrossed in reading}

\label{lesson:RU-C523-zachityvatsya}



\begin{quote}
Before the new one: say the Russian for a heroine, then the Russian for an ending.
\end{quote}

\subsection*{You'll want to know: \ru{зачитываться}}
\textbf{\ru{зачитываться}} (\emph{zachítyvatsya}) — \textquotedblleft{}to get engrossed in reading\textquotedblright{}.

\begin{quote}
\ru{Я} \ru{часто} \ru{зачитываюсь} \ru{до} \ru{глубокой} \ru{ночи}.
\end{quote}

\begin{quote}
\emph{Ya chásto zachítyvayus do glubókoy nóchi.}
\end{quote}

\begin{quote}
I often get so absorbed in a book that I read late into the night.
\end{quote}

Imperfective; the perfective is \ru{зачитаться}. Takes the instrumental: \emph{\ru{зачитываться} \ru{книгой}}.

\begin{cousinweb}
\emph{\ru{зачитываться}} is built on \emph{\ru{за}-} and \emph{\ru{читать}} 'to read', with the reflexive \emph{-\ru{ся}}; EtymDB has no entry.
\end{cousinweb}

\begin{grammarlens}[title={What you've built}]
One more word for telling, explaining and describing at B1.
\end{grammarlens}

\subsection*{Your turn}
\begin{itemize}
\raggedright
  \item \emph{Say it:} \emph{zachítyvatsya}
  \item \emph{Say it:} \emph{zachítyvatsya}, once more
  \item \emph{Say it:} a sentence of your own with \emph{zachítyvatsya}
  \item \emph{Recall:} say the Russian for a heroine, then the Russian for an ending, then say \emph{zachítyvatsya} again
\end{itemize}

\begin{tcolorbox}[breakable,colback=teal!4,colframe=teal!35!black,title={Before you move on}]
What does \ru{зачитываться} mean? (\textquotedblleft{}To get engrossed in reading\textquotedblright{}.) Say it once more.
\end{tcolorbox}
